\documentclass[11pt,a4paper]{article}
\usepackage[utf8]{inputenc}
\usepackage[T1]{fontenc}
\usepackage[spanish,es-noquoting]{babel}
\usepackage[margin=2.4cm]{geometry}
\usepackage{graphicx,booktabs,longtable,array,xcolor,hyperref,fancyhdr,enumitem}
\usepackage{tikz}
\usetikzlibrary{shapes.geometric,arrows.meta,positioning,fit,backgrounds,calc}

\definecolor{itmazul}{RGB}{22,48,110}
\definecolor{acento}{RGB}{0,120,160}
\definecolor{gris}{RGB}{242,244,247}
\definecolor{verde}{RGB}{34,120,60}
\definecolor{rojo}{RGB}{170,40,40}

\hypersetup{colorlinks=true,linkcolor=itmazul,urlcolor=acento}
\pagestyle{fancy}\fancyhf{}
\lhead{\small ComunicaMedellín — Grupo 3}
\rhead{\small Módulo 1 · ITM 2026}
\cfoot{\small\thepage}

\newcommand{\sect}[1]{\textcolor{acento}{\textbf{[Sec.~#1]}}}

\title{\vspace{-1.5cm}
  \textcolor{itmazul}{\Large\textbf{Fan-out Masivo y Presencia en Tiempo Real}}\\[4pt]
  \large ComunicaMedellín — Canal de alertas ciudadanas\\[10pt]
  \normalsize Tópicos Avanzados de Bases de Datos — Módulo 1}
\author{\normalsize Equipo 3}
\date{\normalsize Institución Universitaria ITM · 2026}

\begin{document}
\maketitle
\thispagestyle{fancy}

\begin{center}
\begin{minipage}{0.93\textwidth}\small
\textbf{Resumen.} Diseñamos la arquitectura de un canal de difusión con más de
100.000 suscriptores que debe entregar alertas urgentes de forma casi simultánea
y mantener el estado de presencia de cada usuario. Identificamos cuatro cargas de
trabajo diferenciadas, seleccionamos un motor NoSQL AP con consistencia ajustable
(Apache Cassandra) para la ruta operacional, un log distribuido (Apache Kafka)
para el fan-out particionado y un motor columnar (ClickHouse) para la ruta
analítica. Sostenemos que el sistema debe sacrificar \emph{Consistencia} en
presencia y justificamos esa elección con el impacto de negocio, contrastándola
con el caso TuBoleta, donde la decisión correcta es la opuesta.
\end{minipage}
\end{center}
\vspace{6pt}

% ============================================================
\section{Cargas de trabajo}
% ============================================================
Antes de proponer tecnología identificamos qué hace el sistema. El enunciado
describe cuatro cargas con perfiles distintos \sect{2.1}:

\begin{center}\small
\begin{tabular}{@{}p{0.5cm}p{3.3cm}p{2.4cm}p{4.6cm}p{2.4cm}@{}}
\toprule
\textbf{} & \textbf{Carga} & \textbf{Tipo} & \textbf{Patrón de acceso} & \textbf{Sección} \\
\midrule
C1 & Publicar alerta & OLTP ligero & 1 escritura, baja frecuencia & 2.1 \\
C2 & Fan-out a 100.000 & Streaming & 1 escritura $\rightarrow$ 100K entregas & 4.2 \\
C3 & Historial por usuario & OLTP escritura & 100K inserciones por alerta & 1.3, 2.1 \\
C4 & Presencia (\emph{online}) & Clave-valor & $\sim$10K escrituras/s, alta rotación & 1.2, 1.3 \\
\bottomrule
\end{tabular}
\end{center}

\paragraph{Observación central.} C2 y C4 parecen la misma carga pero exigen
garantías opuestas. La \textbf{presencia} (C4) tolera desactualización: si el
estado de un usuario llega dos segundos tarde, nadie lo nota y ninguna decisión
cambia. La \textbf{entrega de la alerta} (C3) no tolera pérdida: es información
de seguridad pública. Esto obliga a \emph{dos decisiones CAP distintas dentro del
mismo sistema}, y es el eje de nuestro diseño.

\paragraph{Por qué no basta un servidor.} Responde la pregunta (a) del enunciado.
Una arquitectura Cliente-Servidor \sect{3.2} concentra la entrega en un proceso:
100.000 notificaciones salen en serie desde un mismo punto, y el servidor es a la
vez cuello de botella y punto único de falla (SPOF, \sect{3.1}). Alojar ese
servidor en la nube no lo cambia: sigue siendo Cliente-Servidor. Ser distribuido
\sect{3.3} exige repartir datos y procesamiento entre nodos que cooperan mediante
\emph{sharding} y replicación.

% ============================================================
\section{Arquitectura propuesta}
% ============================================================

\begin{figure}[h]\centering
\begin{tikzpicture}[
  font=\small,
  caja/.style={rounded corners=3pt,draw=itmazul!60,fill=gris,
               minimum height=1cm,align=center,inner sep=6pt},
  motor/.style={caja,fill=acento!12,draw=acento,thick},
  flecha/.style={-{Stealth[length=2.5mm]},draw=itmazul!70,thick}]

\node[caja,minimum width=3.2cm] (pub) {Alcaldía\\\scriptsize publica alerta};
\node[caja,minimum width=3.2cm,right=1.4cm of pub] (api) {API\\\scriptsize FastAPI};
\node[motor,minimum width=8.6cm,below=0.9cm of api,xshift=-1.4cm] (kafka)
  {\textbf{Apache Kafka} — log inmutable \sect{4.2}\\[2pt]
   \scriptsize particiones: \texttt{norte} · \texttt{centro} · \texttt{sur} \ \ (sharding por zona, \sect{3.3})};

\node[caja,below=0.9cm of kafka,xshift=-3cm,minimum width=2.1cm] (w1) {worker 1};
\node[caja,below=0.9cm of kafka,minimum width=2.1cm] (w2) {worker 2};
\node[caja,below=0.9cm of kafka,xshift=3cm,minimum width=2.1cm] (w3) {worker 3};

\node[motor,below=1cm of w2,xshift=-2.6cm,minimum width=5cm] (cass)
  {\textbf{Cassandra} · 3 nodos, RF=3\\
   \scriptsize AP · \texttt{ONE}/\texttt{QUORUM} \sect{1.3, 3.3}\\
   \scriptsize presencia + historial};
\node[motor,below=1cm of w2,xshift=3.1cm,minimum width=4.2cm] (ch)
  {\textbf{ClickHouse}\\\scriptsize columnar OLAP \sect{2.2, 2.4}\\
   \scriptsize métricas agregadas};

\draw[flecha] (pub) -- (api);
\draw[flecha] (api) -- (kafka);
\draw[flecha] (kafka.south) -- (w1.north);
\draw[flecha] (kafka.south) -- (w2.north);
\draw[flecha] (kafka.south) -- (w3.north);
\draw[flecha] (w1.south) -- (cass.north west);
\draw[flecha] (w2.south) -- (cass.north);
\draw[flecha] (w3.south) -- ($(cass.north east)+(0,0)$);
\draw[flecha,dashed,draw=verde!70] (kafka.east) -- ++(1.1,0) |- (ch.east)
  node[midway,right,align=left,font=\scriptsize,text=verde!60!black,xshift=2pt]
  {Kafka\\Engine};
\end{tikzpicture}
\caption{Arquitectura distribuida. La ruta operacional (Cassandra) y la analítica
(ClickHouse) se alimentan del mismo log, configurando una arquitectura
Lambda \sect{4.1}.}
\end{figure}

\subsection{Justificación de cada decisión}

\begin{description}[leftmargin=0.9cm,style=nextline]
\item[Apache Cassandra \sect{1.3, 3.3}]
El módulo la clasifica como \emph{«Cassandra (AP, wide-column)»} y documenta el
caso Netflix: arquitectura \emph{peer-to-peer} sin punto único de falla, más de un
billón de peticiones diarias. Es el único motor del stack de referencia que ofrece
\textbf{consistencia ajustable por consulta}, requisito literal del reto
(\texttt{ONE} frente a \texttt{QUORUM}).

\item[Modelo clave-valor \sect{1.3}]
Responde la pregunta (b). El patrón de acceso a presencia es único: obtener el
estado de un usuario por su identificador. No hay \textsc{join}s, ni integridad
referencial que preservar, ni consultas por atributos secundarios. Un modelo
normalizado 3NF \sect{2.1} añadiría costo sin aportar garantías útiles.

\item[Apache Kafka \sect{4.2}]
Log distribuido, inmutable y particionado. Responde la pregunta (d): la zona actúa
como clave de partición, de modo que cada worker consume una partición
independiente. Escalar consiste en añadir particiones y consumidores, no en
agrandar un servidor.

\item[ClickHouse \sect{2.2, 2.4}]
Cassandra está optimizada para lectura por clave de partición; una agregación
\texttt{GROUP BY} sobre 100.000 usuarios la obligaría a recorrer el clúster
completo. Es el conflicto OLTP/OLAP de la \sect{2.2}. ClickHouse resuelve esas
agregaciones en milisegundos, como en el caso Mercado Libre \sect{2.4}.

\item[Alternativa descartada: CockroachDB \sect{1.3, 3.3}]
Es la opción correcta para TuBoleta, no para este caso. Al ser CP con consenso
Raft, garantiza serializabilidad en toda escritura: los nodos nunca divergen. Eso
hace que el trade-off que debemos demostrar sea inobservable, y encarece cada
\emph{heartbeat} de presencia sin beneficio.
\end{description}

% ============================================================
\section{Análisis CAP}
% ============================================================
El teorema \sect{1.2} establece que ante una partición de red sólo pueden
sostenerse dos de las tres propiedades. Como las particiones son inevitables en
un sistema distribuido, la elección real es entre CP y AP.

\begin{figure}[h]\centering
\begin{tikzpicture}[font=\small]
  \def\R{2.05}
  \coordinate (C) at (90:\R);
  \coordinate (A) at (210:\R);
  \coordinate (P) at (330:\R);
  \fill[acento!8] (C) -- (A) -- (P) -- cycle;
  \draw[thick,draw=itmazul!70] (C) -- (A) -- (P) -- cycle;
  \node[above=2pt] at (C) {\textbf{C}onsistencia};
  \node[below left=1pt] at (A) {\textbf{A}vailability};
  \node[below right=1pt] at (P) {\textbf{P}artition tol.};
  \draw[very thick,rojo] ($(C)!0.5!(A)+(-0.15,0.15)$) -- ++(0.5,-0.5);
  \node[align=center,font=\scriptsize] at (0,-0.35)
    {\textbf{AP}\\ presencia\\ Cassandra};
  \node[right=2.6cm of P,align=left,font=\scriptsize,anchor=west,yshift=1.4cm]
    {\textcolor{rojo}{$\times$} \textbf{Se sacrifica C} en presencia.\\[4pt]
     Un desfase de 2\,s en «en línea»\\ no altera ninguna decisión\\ de negocio.};
\end{tikzpicture}
\caption{Elección CAP para la carga de presencia (C4).}
\end{figure}

\subsection{Dos decisiones, no una}

\noindent\textbf{Presencia (C4) $\rightarrow$ AP, nivel \texttt{ONE}.}
Se sacrifica Consistencia. \emph{Impacto de negocio:} nulo. Nadie toma una
decisión a partir del estado «en línea» de un tercero. En cambio, elegir CP sería
costoso: durante una partición el sistema rechazaría \emph{heartbeats}, la
aplicación mostraría a todos los usuarios desconectados y la percepción sería la
de un servicio caído — peor que un dato con dos segundos de antigüedad.

\noindent\textbf{Historial de alertas (C3) $\rightarrow$ AP, escritura \texttt{QUORUM}.}
La alerta no puede perderse: comunica cierres viales y cancelaciones de clase.
Escribir con \texttt{QUORUM} (2 de 3 réplicas) asegura durabilidad ante la caída
de un nodo sin renunciar a la disponibilidad. El sobrecosto ($\sim$40\,ms) es
irrelevante porque la publicación es poco frecuente.

\subsection{Contraste con TuBoleta}

\begin{center}\small
\begin{tabular}{@{}p{3.1cm}p{5.2cm}p{5.2cm}@{}}
\toprule
 & \textbf{TuBoleta (Equipo 1)} & \textbf{ComunicaMedellín (Equipo 3)} \\
\midrule
Recurso & Silla A-12: finito y exclusivo & Mensaje: replicable sin límite \\
Conflicto & Dos compradores compiten & Nadie compite \\
Si falla C & Doble venta, contingencia legal & Presencia desfasada 2\,s \\
Decisión & \textbf{CP} (Raft, serializable) & \textbf{AP} (eventual, ajustable) \\
Motor & CockroachDB & Cassandra \\
\bottomrule
\end{tabular}
\end{center}

\noindent La consistencia fuerte cuesta latencia, y ese costo se justifica cuando
existe un recurso en disputa. En presencia no hay disputa: nadie «gana» el estar
en línea. Pagar consenso por cada actualización sería gasto sin contrapartida.

% ============================================================
\section{Modelo de datos}
% ============================================================
\begin{verbatim}
CREATE KEYSPACE comunica WITH replication =
  {'class':'SimpleStrategy','replication_factor':3};

CREATE TABLE presencia (            -- C4: clave-valor puro
    usuario_id  text PRIMARY KEY,   -- clave de partición: reparte por hash
    estado      text,               -- 'online' | 'offline'
    last_seen   timestamp,
    zona        text,
    dispositivo text
);

CREATE TABLE historial_usuario (    -- C3: fan-out on write
    usuario_id text,                -- partición: historial junto por usuario
    enviado_en timestamp,           -- clustering: recientes primero
    alerta_id  uuid,
    titulo     text,
    severidad  text,
    PRIMARY KEY ((usuario_id), enviado_en)
) WITH CLUSTERING ORDER BY (enviado_en DESC);
\end{verbatim}

\noindent\textbf{Estrategia \emph{fan-out on write}.} Al publicar una alerta se
escribe una fila por suscriptor. Encarece la escritura (que ocurre una vez) para
abaratar la lectura (que ocurre 100.000 veces): cada usuario abre su historial
leyendo una sola partición. Es el intercambio correcto dado el patrón de acceso.

% ============================================================
\section{Reto práctico}
% ============================================================
La demostración se ejecuta con \texttt{make}. Cada bloque cubre un punto del
enunciado:

\begin{center}\small
\begin{tabular}{@{}p{4.6cm}p{6.2cm}p{3.4cm}@{}}
\toprule
\textbf{Comando} & \textbf{Qué demuestra} & \textbf{Requisito} \\
\midrule
\texttt{make up} & Clúster de 3 nodos operativo & Exigido \\
\texttt{make seed} & 10.000 usuarios concurrentes & Exigido \\
\texttt{make demo-consistencia} & \texttt{ONE} frente a \texttt{QUORUM}, dos nodos & Exigido \\
\texttt{make demo-fanout} & Fan-out particionado por zona & Exigido \\
\texttt{make demo-nodo-caido} & Disponibilidad con un nodo caído & \emph{Extra} \\
\texttt{make ch-demo} & Agregaciones columnares & \emph{Extra} \\
\texttt{make kafka-replay} & Relectura desde \emph{offset} 0 (Kappa) & \emph{Extra} \\
\texttt{make zk-shell} & Coordinación con Zookeeper & \emph{Extra} \\
\bottomrule
\end{tabular}
\end{center}

\subsection{Metodología: por qué se simula latencia de red}
Los tres nodos comparten una máquina y replican en microsegundos. Sin latencia,
\texttt{ONE} nunca devuelve un dato desactualizado y la demostración no probaría
nada: mediría el laboratorio, no el sistema. Se introducen 250\,ms con
\texttt{tc netem} de forma \textbf{asimétrica} --- sólo al tráfico de replicación
entre nodos (puerto 7000), no a las lecturas del cliente (puerto 9042) --- de modo
que el cliente lea de una réplica que aún no recibió la escritura. En producción
esa latencia existe siempre; el laboratorio la elimina y aquí se restituye.

\subsection{Resultado principal: el teorema CAP, medido}
Escritura confirmada con \texttt{ONE} y lectura inmediata y simultánea desde el
nodo 1 y el nodo 3, sobre 150 rondas:

\begin{center}\small
\begin{tabular}{@{}lrrl@{}}
\toprule
\textbf{Nivel} & \textbf{Divergencias} & \textbf{Latencia P50} & \textbf{Lectura} \\
\midrule
\texttt{ONE}    & \textcolor{rojo}{\textbf{149/150 (99{,}3\,\%)}} & \textbf{1{,}14\,ms}
  & Rápido, puede ir atrasado \\
\texttt{QUORUM} & \textcolor{verde}{\textbf{0/150 (0{,}0\,\%)}} & \textbf{130{,}3\,ms}
  & Consistente, 114$\times$ más lento \\
\bottomrule
\end{tabular}
\end{center}

\noindent\textbf{\texttt{QUORUM} elimina la inconsistencia por completo, a costa de
ser 114 veces más lento.} Se elige \texttt{ONE} para presencia porque la
inconsistencia no tiene costo de negocio y la latencia sí.

\subsection{Resultados complementarios}
\begin{center}\small
\begin{tabular}{@{}p{4.4cm}p{4.4cm}p{5.2cm}@{}}
\toprule
\textbf{Demostración} & \textbf{Resultado medido} & \textbf{Qué prueba} \\
\midrule
Carga de suscriptores & 10.000 en 1{,}3\,s (7.597 escrituras/s)
  & El clúster absorbe escritura concurrente \\
Fan-out particionado & 10.127 entregas/s; 2{,}6$\times$ frente a secuencial
  & El paralelismo por zona es real \\
Reparto por zona & 34 / 33 / 33
  & Una zona por partición y worker \\
Caída de un nodo & \texttt{ONE} 3{,}1\,ms; \texttt{QUORUM} 752\,ms; \texttt{ALL} falla
  & A mayor consistencia exigida, menor disponibilidad \\
Compresión columnar & 1 millón de filas, ratio 2{,}4:1
  & ClickHouse lee sólo las columnas necesarias \\
\bottomrule
\end{tabular}
\end{center}

\noindent La interfaz web (\texttt{make api} y \texttt{make frontend}) permite
publicar una alerta como funcionario de la Alcaldía y observar el recorrido real
--- Kafka, los tres \emph{workers}, Cassandra y los dispositivos --- además de
comparar la lectura del mismo usuario desde los tres nodos con un conmutador
\texttt{ONE}/\texttt{QUORUM} en caliente.

% ============================================================
\section{Conclusiones}
% ============================================================
\begin{enumerate}[leftmargin=*,itemsep=3pt]
\item El caso contiene cuatro cargas y dos decisiones CAP distintas. Tratarlo con
una única política habría sido incorrecto.
\item La consistencia es un parámetro por consulta, no una propiedad global del
sistema: \texttt{ONE} para presencia, \texttt{QUORUM} para entrega de alertas.
\item El sacrificio de C se justifica por el impacto de negocio, no por
conveniencia técnica: en presencia ese impacto es nulo, mientras que sacrificar A
degradaría el canal justo cuando más se necesita.
\item Separar la ruta operacional (Cassandra) de la analítica (ClickHouse) evita
el conflicto OLTP/OLAP descrito en la \sect{2.2}.
\end{enumerate}

\vfill
\begin{center}\footnotesize\color{gray}
Repositorio: \texttt{github.com/lina082/comunicamedellin-grupo3} ·
Todo el stack es software libre (Apache 2.0 / MIT).
\end{center}

\end{document}
